\properlane{deliverance-and-waiting}{Deliverance begun, and the people who still wait}

Paul thanks the Father for having rescued the faithful from darkness.
He also asks for strength to endure. The rescue is an accomplished
benefit, while the endurance is still required. Paul's thanksgiving
concerns a gift already received; his petition concerns the faithful's
unfinished growth.

The distinction is between redemption already received and the
deliverance still awaited in resurrection and final gathering.
Augustine develops that distinction through the Psalms. Chrysostom
explains why Paul's continuing prayer and Christ's warnings are
addressed to people God has already helped. Read together, their
arguments permit confidence through affliction without making
affliction evidence that Christ's promise has failed.

\subsection{Why the Gradual can speak as though rescue were past}

The Gradual says that God has delivered his people and put their enemies
to shame. The full Psalm then laments that they have been defeated and
scattered. Its speakers have not simply forgotten the earlier praise.
They ask why their present condition differs so sharply from what they
have heard about God and from what they hope to receive.

Augustine attends to the verbs. After speaking of rescue as something
done, the psalmist says that the people will praise God. For Augustine,
those future verbs show that the rescue described in the past tense is
itself still to come. He compares Psalm 21(22), which speaks of pierced
hands and feet as an event already accomplished even though Christ's
Passion was still future for the prophet. What God knows will happen
can be spoken of with the certainty of what has
happened.\footnote{Augustine, \work{Enarrationes in Psalmos},
Ps 43, Latin §§8--9 (PL 36) and NPNF Ps XLIV §§6--7.}

His next step concerns the present lament. The ancestors' deliverance
is past, the final kingdom is future, and the Church suffers between
them. God has not become another God during that interval.
Augustine's explanation therefore preserves three things that a
simple victory song might leave apart: remembered action, promised
fulfilment and present pain. The Psalm's question is intelligible
precisely because the promise is still believed.\footnote{Augustine,
\work{Enarrationes in Psalmos}, NPNF Ps XLIV §§4, 8, 13.}

Bellarmine gives another account of the appointed past tense. He takes
it as thanksgiving for rescues actually received. The people then
describe a later affliction. On this account confidence rests on
God's previous action rather than on a prophetic announcement phrased
as completed action. The difference concerns how the verse says what
it says. Both interpretations recognize that the faithful may praise
deliverance while they suffer; only Augustine's makes this particular
perfect tense refer prophetically to the future.\footnote{Bellarmine,
\work{Commentary on the Book of Psalms}, Psalm 43, explanation of
native vv. 7--8 (Vulgate vv. 8--9), then native v. 9 for the onset
of the lament.}

For this Mass, the consequence is that thanksgiving need not be
measured by an absence of danger. The Gradual's praise is followed
by the Alleluia's cry from the depths. A person who understands the
Psalm's complete argument can sing both. The cry acknowledges a
present need; the thanksgiving rests on an action or promise of
God whose reach exceeds the present circumstances.

\subsection{The exiles' peace includes a long obedience}

The Introit's promise also concerns a people whose situation is not
immediately reversed. Jeremiah opposes the expectation of a quick
return from Babylon. The letter's instructions to build and plant
would be cruelly misleading if God's peace meant that they need
only wait briefly before abandoning that life. Instead, faithful
work and prayer for the city become possible within the
waiting.\footnote{Jer 29:4--14.}

Jerome takes the promised return as future and opposes the false
prophets' assurances of an early release. In his exposition of the
first part of the promise, he also reads the waiting as looking ahead
to rest after the troubles of this life. In the next passage, which
begins with the gathering clause used in the chant, he explains the
gathering of the historical exiles on the promised
timetable.\footnote{Jerome, \work{In Ieremiam} V,
\textit{Vers. 10 seqq.} (Jer 29:10--14a) and the opening of
\textit{Vers. 14 seqq.} (Jer 29:14b); PL 24, cols. 893--894
(facsimile PDF p. 451).}

The Introit's psalm verse puts the end of captivity in the past tense,
while the whole Psalm still asks for conversion. Augustine again hears
prophetic certainty. He reasons from the later petition: if the restoration
had already happened at the prophet's standpoint, why would the prophet
still ask for it? The past form announces a future gift. His explanation
of captivity also reaches beyond external displacement to sin and
mortality.\footnote{Augustine, \work{Enarrationes} 84.1--13,
especially §§2--5.}

For Christians giving thanks with Paul for redemption already received,
this account of captivity permits a further application. Forgiveness has
begun their release, while freedom from mortality remains ahead. That
application follows from setting the Psalm beside the Epistle's rescue
and continuing prayer. In the Collect, too, those called faithful ask
for wills renewed in fruitful work and for further remedies. Fromage
explains their good works as grace-assisted preparation for greater
grace.\footnote{Col 1:9--14; Fromage, vol. XI, p. 503.}

This is a specific kind of patience. Jeremiah's people are to live
responsibly in a place from which they hope to depart. The
worshipper can likewise pursue the fruit named by the Collect
without supposing that earthly circumstances must first become
secure. Hope for restoration permits sustained work during
uncertainty.

\subsection{What has already happened in the Son's kingdom}

Chrysostom explains Paul's rescue and transfer as two benefits.
Freedom from an oppressor is already a gift. Admission into the
beloved Son's kingdom adds a positive belonging and inheritance.
The Father has also fitted the recipients to share what he gives.
Paul's thanksgiving concerns this actual change in their condition
before God.\footnote{Chrysostom, \work{Homilies on Colossians}
2, on Col 1:12--14.}

Yet Paul has not ceased praying. Chrysostom returns to the repeated
requests for fullness: the Colossians have received something real,
but they still need to grow. Their existing faith and love are
reasons to encourage them, much as competitors near victory receive
encouragement to continue. A gift can be gratefully acknowledged
without describing its recipient as already perfected.

The requested strength includes endurance. Aquinas distinguishes
failure caused by the difficulty of adversity from failure caused
by the delay of the promised reward. Patience answers the first;
longanimity answers the second. His distinction explains why delay
can become a trial even when no new disaster occurs. The good is
still absent, and the person must continue wanting and pursuing
it.\footnote{Aquinas, \work{Super Colossenses} 1, lectio 3,
on Col 1:11. Chrysostom, homily 2, also distinguishes patience
and long-suffering, but uses the contrast of those whom one can
and cannot requite.}

The Epistle requests joy alongside these virtues. Such joy
concerns the received gift and its promised end. It need not
erase grief at what is suffered. Chrysostom himself uses Job's
thanksgiving to distinguish gratitude from a claim that loss
caused no anguish. A suffering Christian may therefore give
thanks truthfully without denying the suffering.\footnote{Chrysostom,
\work{Homilies on Colossians} 2, on the transition from 1:11 to 1:12.}

\subsection{The warning to flee and the promise to gather}

The Gospel presents a harder difficulty than delayed consolation.
Jesus warns of an emergency in which a person must leave the
house without collecting possessions. Winter and the Sabbath add
obstacles to flight from Judea. Chrysostom asks how a pregnant woman
could make the necessary escape and how a nursing mother could leave
her infant behind.\footnote{Chrysostom,
\work{Homilies on Matthew} 76, on Mt 24:16--22.}

The command to flee also requires a response. Trusting Christ
includes heeding his warning and taking the action he names.
The shortening of the days for the elect expresses God's care
within the danger. It does not promise that the danger will
never begin, or that each believer will be exempt from death.

Chrysostom then distinguishes the warning about Jerusalem from
the account of Christ's final coming. The word ``then'' at
verse 23 need not mean the next moment after the siege. He
compares Scripture's use of a similar transition between the
childhood narrative and John the Baptist's ministry. The
transition introduces the period being discussed without
measuring the intervening interval.\footnote{Chrysostom,
\work{Homilies on Matthew} 76.2; Aquinas,
\work{Super Matthaeum} 24, lectio 3, makes the same temporal
distinction.}

The word ``immediately'' in verse 29 raises a further question.
Chrysostom does not treat it as another indefinite ``then.'' He refers
the immediately preceding tribulation to Antichrist and the false prophets,
whose deception he has just explained. The public coming follows that
final trial promptly. His distinction therefore depends on which
tribulation Jesus is discussing, rather than on stretching every temporal
word. The Judean emergency and the last deception remain different
afflictions, each limited for the elect's sake.\footnote{Chrysostom,
\work{Homilies on Matthew} 76.2--3, on Mt 24:23--29.}

This explanation preserves both the local urgency of flight and
the universal manifestation of the Son of man. A reading confined
to one siege readily explains the roof, field and season.
Chrysostom judges that the public coming and gathering reach
further. Conversely, an interpretation concerned only with a
future catastrophe would leave the Judean directions without
their concrete force.

The false claimants promise a Christ who must be found in a
particular hidden place. Lightning requires no privileged
informant to announce where it has been concealed. Christ's
coming will be manifest, and his angels will gather the elect
from the four winds. The distinction is about recognition:
the faithful need not entrust their hope to someone offering
secret access to the returning Lord.

The saying about the body and the eagles adds an image of gathering.
Chrysostom identifies the gathered multitude with angels, martyrs and
saints. Aquinas relates the body to Christ's Passion and the eagles to
the saints drawn to him. The image thus answers the question of where
Christ's people will assemble: they gather around Christ himself.
The later promise that angels will gather the elect develops that
assurance. Their destination is determined by him, without a secret
claimant supplying directions.\footnote{Chrysostom,
\work{Homilies on Matthew} 76.3; Aquinas, \work{Super Matthaeum}
24, lectio 3, on Mt 24:28, 30--31.}

The difficult phrase ``this generation'' receives a related
explanation in Chrysostom. He understands it as the generation
of believers, identified by faith and practice rather than
only by a shared lifespan. The disasters will not extinguish
this people. The promise concerns the Church's endurance
through what Jesus has described. It supplies no guarantee
that the individuals first hearing the discourse will all
remain alive until its final fulfilment.\footnote{Chrysostom,
\work{Homilies on Matthew} 77, on Mt 24:33--35; compare
Aquinas, \work{Super Matthaeum} 24, lectio 3, on 24:34.}

\subsection{Hope before sight, and medicine before resurrection}

Augustine's exposition of Psalm 129 makes the remaining interval
concrete through Christ's risen body. In the whole Psalm's
morning watch he hears Christ's Resurrection. Sins have been
forgiven, but the faithful have not yet risen as Christ has
risen. What has happened in the Head is the reason the members
hope for the same fulfilment. His argument moves from Christ's
assumption of our flesh, through its sacrifice, to its rising
as the first fruits of ours.\footnote{Augustine,
\work{Enarrationes} 129.3--8. The morning watch occurs in
the complete Psalm, beyond the two verses appointed in this Mass.}

The Offertory's repeated cry can be understood within that
waiting. People who hope for resurrection still bring their
need before God. The Secret asks them to pass from earthly
cravings to heavenly desires. Here that change includes
ceasing to treat survival, possessions or an undisturbed
present as the final fulfilment of God's peace.

In the Trinity Preface the Church gives thanks to the Father
``at all times, and in all places'' and confesses the Father,
Son and Holy Spirit in one essence and equal majesty. The
Preface says that the angels praise this same God day by day
``with one voice,'' and their words follow in the Sanctus.
When the congregation joins in ``Holy, holy, holy, Lord God
of hosts,'' it shares already in the praise it hopes to offer
without end.\footnote{1962 Missal, p. 293, no. 1082;
Lasance/Walsh, pp. 773--774.}

At Communion, Christ's assurance permits prayer before the
answer is visible. Its qualification by faith and forgiveness
will receive fuller explanation in the next interpretation.
Here the Postcommunion makes the continuing need explicit:
the received mysteries are asked to cure remaining faults.
Augustine describes the Eucharistic food and drink as received
for spiritual health through God's invisible operation.
Fromage explains this particular closing petition as the
cure of the communicant's remaining defects.\footnote{Bede,
\work{In Marci evangelium expositio} III, on Mk 11:22--23, 25;
Augustine, \work{De Trinitate} III.4.10; Fromage, vol. XI,
p. 513.}

The final healing for which this interpretation hopes includes
resurrection. Until then, endurance, truthful discernment,
fruitful work and repeated petition remain possible responses
to grace. The gathered people depend on the Christ whose words
abide when the visible world changes.

\begin{fourSenses}
\item[Literal.] Jeremiah promises return to the deportees. The Gradual
speaks of rescue in the past tense, which Augustine reads as prophetic
and Bellarmine as remembered. The cry from the depths voices a present
need. Paul prays for believers already rescued, and Jesus warns of danger
and announces his public coming. The prayers ask further grace and cure.
\item[Allegorical.] The Father brings the redeemed into the
beloved Son's kingdom through his Blood. Christ's risen flesh
is the first fruits of the resurrection his members await.
The spiritual fulfilment of captivity and return retains
Jeremiah's historical exiles as its first addressees.
\item[Moral.] Persevere during adversity and delay, heed
Christ's warnings, refuse secret claimants and continue the
good work for which the Collect asks renewed wills.
\item[Anagogical.] The final gathering and resurrection
complete the deliverance already begun. Thanksgiving with
the angels anticipates unending praise of the triune God.
\end{fourSenses}
